\section{Computation in FEVC}\label{sec:computation}

In large samples FEVC computes $\hat\theta$ without forming an inverse. It solves the normal equations for
$\hat\beta$, and it obtains the leverages $P_{ii}$ and the weights $B_{ii}$ from random projections, the
Johnson--Lindenstrauss approximation (JLA) of KSS: each projection is a solve of the normal equations with a
random right-hand side. The standard error reuses these quantities and adds a few solves of the same kind.

\begin{enumerate}
\item \emph{Variance estimates.} The $q_g$ take one pass over the stayers' rows. The model for movers takes two
least-squares fits with five and six coefficients; their inputs ($\hat\alpha$, $\hat\psi$, $P_{ii}$,
$\hat\sigma_i^2$) are by-products of the point estimate.
\item \emph{The influences $\hat m$.} With $d_i=B_{ii}/(1-P_{ii})$ and $D=\diag(d_i)$, the weights are
$C=B-\frac12(DM+MD)$, and
\[
\hat m=Cy=X\,S_{xx}^{-1}\big(A\hat\beta+\tfrac12X'Dy\big)-\tfrac12D\,(y-X\hat\beta)-\tfrac12Dy ,
\]
which is one solve per component.
\item \emph{The double sum.} It is not computed pair by pair. Draw $D_s$ vectors $\zeta^{(d)}$ with
independent normal entries of mean zero and variances $\st_i^2$, apply the weights to each as in step~2, and
use
\begin{equation}\label{eq:sim}
\sum_i\sum_{\ell\ne i}C_{i\ell}^2\,\st_i^2\,\st_\ell^2\;\approx\;
\frac{1}{D_s}\sum_{d=1}^{D_s}\sum_i\st_i^2\,\big(C\zeta^{(d)}\big)_i^{2} ,
\end{equation}
which is right on average because the expectation of $(C\zeta)_i^2$ over the draws is
$\sum_\ell C_{i\ell}^2\st_\ell^2$. Each draw costs one solve per component and one for its residual. The same
draws serve all components and their covariances.
\item \emph{The adjustment} is the closed form in \eqref{eq:Vhat}.
\end{enumerate}

\begin{proposition}[Simulated double sum]\label{prop:sim}
Let $\hat\V_{D_s}$ be \eqref{eq:Vhat} with the double sum computed by \eqref{eq:sim}. Under the assumptions of
Proposition~\ref{prop:conc}, with $0\le\st_i^2\le\bth\sigma_i^2$ on movers' rows,
$\E[\hat\V_{D_s}\mid\text{data}]=\hat\V[\hat\theta]$ and
\[
\frac{\E\big[(\hat\V_{D_s}-\hat\V[\hat\theta])^2\big]}{\V[\hat\theta]^2}\;\le\;\frac{c''\,\bth^4}{D_s}\cdot\frac{\lambda_{\max}^2}{\V[\hat\theta]} ,
\]
where $c''$ depends only on $\bar c$.
\end{proposition}

The noise from the draws is thus small under the same condition (b) that makes the variance estimate settle
down, and it falls with the number of draws. (The bound is for non-random values on movers' rows. With fitted
values the first statement still holds, and given the data the variance of $\hat\V_{D_s}$ is
$8\tr(\tilde K^4)/D_s$ with $\tilde K_{i\ell}=\st_iC_{i\ell}\st_\ell$.) We use $D_s=16$; Section~\ref{sec:evidence} reports how much two
runs differ.

\paragraph{Stayers' leverages are known exactly.} A stayer's rows move only that stayer's person effect, so
\[
P_{ii}=\frac{1}{T_g},\qquad B_{ii}=\frac{A_{gg}}{T_g^{\,2}}
\]
for every row of stayer $g$, and \eqref{eq:cg} follows. The prototype uses these values instead of projected
ones for stayers. This costs nothing and removes the projection noise from three quarters of the rows of the
panel in Section~\ref{sec:evidence}.

\paragraph{Projected leverages for movers.} Movers' $P_{ii}$ and $B_{ii}$ come from projections, so the weights
actually used have a small non-zero diagonal. The propositions are for exact weights; every simulation below
uses the projected ones, so the evidence covers them. One point concerns the estimate rather than its standard
error. The projection error of $\hat\theta$ is not sampling error and is not part of the standard error, and in a
very large panel the two can be of the same size. On the Veneto panel the standard error of
$\hat\sigma^2_{\alpha+\psi}$ is \R{real.rel_se_sum} of the estimate; with 200 projections and no further
correction the prototype's estimate was off by up to \R{sim.bias200} standard errors. With 1{,}600 the bias is
smaller but still systematic: up to about 0.2 standard errors, and downwards in every scenario. FEVC applies an analytical correction for the finite number of projections; the prototype
does not, and uses 1{,}600 projections throughout this note.

\paragraph{Cost.} On the six-year Veneto panel (4.1 million rows, 87{,}000 firms) the prototype factors the
firm system once, which took between 8 and 37 minutes depending on the machine. A standard error for all four
components then needs 67 solves beyond the least-squares fit: three for $\hat m$ and four for each of the 16
draws (one for the residuals of the draw and one for each of three components; the weights for the variance of
the sum are a linear combination of the other three). In a run of this model alone the least-squares fit with the three influence solves took
\R{real8.seconds.fit} seconds and the 64 solves of the draws \R{real8.seconds.draws} seconds. This compares with
\R{real.probe_solves} solves for the 1{,}600 projections of the point estimate. FEVC itself solves the
normal equations iteratively, without a factorisation, and its point estimate takes roughly 600 to 800 solves;
the 67 solves of the standard error would add about a tenth to that. This is an expectation from the solve
counts, not a measurement: the estimator is not yet part of FEVC. Appendix~\ref{app:cost} explains why so few
solves are enough, by comparison with two earlier computations.
